\clearpage
\fancyhead[L]{\small DELAMBRE. TABLES ÉCLIPTIQUES.}
\hypertarget{native-608-lower}{}
\begin{center}\small Göttingische gelehrte Anzeigen. Stück 45. Seiten 441..443. 19. März 1818.\\[3pt]\rule{.72\linewidth}{.3pt}\end{center}

\emph{Tables écliptiques des satellites de jupiter, d'après la théorie de M. le marquis de \textup{\textsc{Laplace}}, et la totalité des observations faites depuis \textup{1662} jusqu'à l'an \textup{1802}; par M. \textup{\textsc{Delambre}}, chevalier de St. Michel etc.} Paris 1817. Bei der Witwe Courcier. Die Einleitung 58 Seiten, die Tafeln 32 halbe Bogen. 4.

\begingroup\brokenpenalty=10000\relax
Die vor beinahe 30 Jahren von Herrn \textsc{Delambre} nach \textsc{Laplace}s Theorie berechneten und in der dritten Ausgabe von \textsc{Lalande}s Astronomie abgedruckten Tafeln für die Verfinsterungen der Jupitertrabanten wurden bekanntlich von den Astronomen mit verdientem Beifall aufgenommen und sind in allgemeinen Gebrauch gekommen. Jener arbeitsame und unverdrossene Rechner hat seitdem die zahlreicher gewordenen Beobachtungen und die noch größere Vervollkommnung der Theorie genutzt, um die Grundlagen seiner Tafeln von Neuem zu verbessern. So sind die vorliegenden neuen Tafeln entstanden, die bereits vor mehreren Jahren gedruckt, aber erst jetzt ausgegeben worden sind. Der große Geometer, unter dessen Leitung diese mühsame Arbeit ausgeführt worden ist, hat ihr bereits das ehrenvolle Zeugnis gegeben: `\textsc{Delambre} a exécuté ce travail important avec le plus grand succès; et ses tables qui réprésentent les observations avec l'exactitude des observations mêmes offrent au navigateur un moyen sûr et facile pour avoir sur le champ, par les eclipses des satellites, et surtout par celles du premier, la longitude des lieux où il atterre.' Die Einrichtung der Tafeln unterscheidet sich wenig von derjenigen der früheren, und die Einleitung gibt eine sehr ausführliche Anleitung zu ihrem Gebrauch. Nur hätten wir gewünscht, dass der Verfasser die Ergebnisse seiner Verbesserungsarbeit ausführlicher mitgeteilt hätte, sodass man die Änderungen, welche die einzelnen durch die Beobachtungen zu bestimmenden konstanten Grundlagen (deren Zahl sich, die Lichtgeschwindigkeit mitgerechnet, be\SourceLeaf{609}kanntlich auf zweiunddreißig beläuft) erfahren haben, klar und leicht hätte überblicken und den Grad der ihnen beizulegenden Genauigkeit hätte beurteilen können. Aus einer Stelle der Einleitung lässt sich schließen, dass alle diese Änderungen sehr gering gewesen sind. Er erklärt, dass, obwohl die Fehler der neuen Tafeln im Allgemeinen geringer sind als die der früheren, `ce que nous avons gagné ne vaut peut-être
pas le travail qu'il a couté', ganz im Geiste der Ansicht, den Wert einer solchen Arbeit allein nach ihrem praktischen Nutzen für die Bestimmung der geografischen Länge zu bemessen. Nur zum Koeffizienten der Lichtgleichung, den Herr \textsc{Delambre} mit $493''2$ ansetzt, bemerkt er, dass sich diese Bestimmung auf mehr als 1000 Beobachtungen des ersten Trabanten gründe und dass die davon nicht merklich verschiedene frühere Bestimmung auf 500 Beobachtungen beruht habe. Es ist doch sehr bemerkenswert, dass dieser Koeffizient die Aberrationskonstante $20''25$ ergibt, während diese, aus den Rektaszensionen des Polarsterns bestimmt, etwas größer ausfällt. Es wäre indessen vielleicht voreilig, hieraus schon auf einen Unterschied zwischen der Geschwindigkeit des Fixsternlichts und der des Trabantenlichts oder zwischen der Lichtgeschwindigkeit am Umfang der Erdbahn und der näher zur Sonne hin schließen zu wollen. Doch verdient dieser Umstand gewiss fortgesetzte Aufmerksamkeit.
\par\clearpage\endgroup
